\documentclass[a4paper,11pt,twoside]{article}
\usepackage[utf8]{inputenc}
\usepackage[left=2.5cm,right=2.5cm, top=3cm, bottom=3cm]{geometry}
\usepackage[dvipsnames]{xcolor}
%\usepackage{xcolor,colortbl}
%\documentclass[xcolor=dvipsnames]{beamer} 
\usepackage{graphicx}
\usepackage{multirow}
\usepackage{enumitem}
%\usepackage[demo]{graphicx}
\usepackage{caption}
\usepackage{subcaption}
%\usepackage{subfig}
\usepackage{graphicx, caption, subcaption}
\usepackage{float}
\usepackage{amsmath}
%\usepackage[english]{babel}
%\usepackage{blindtext}
\usepackage{titlesec}

% for braket notation
\usepackage{mathtools}
\DeclarePairedDelimiter\bra{\langle}{\rvert}
\DeclarePairedDelimiter\ket{\lvert}{\rangle}
\DeclarePairedDelimiterX\braket[2]{\langle}{\rangle}{#1\,\delimsize\vert\,\mathopen{}#2}

%bibliographie
\usepackage{hyperref}
\usepackage{xcolor}

%\usepackage{biblatex}
%\addbibresource{main.bib}
%inhaltsverzeichnis
\setcounter{tocdepth}{2}
\titleformat*{\section}{\huge\bfseries}
\titleformat*{\subsection}{\Large\bfseries}
\titleformat*{\subsubsection}{\large\bfseries}

%Kopf und Fusszeile
\usepackage{fancyhdr}
\pagestyle{fancy}
\setlength{\headheight}{14pt} % Fixes the "headheight too small" warning
\renewcommand{\sectionmark}[1]{\markboth{#1}{}} 
\fancyhf{}
\fancyhead[R]{\textit{Illia Yakushevskyi}} % Put your actual name here
\fancyhead[L]{\leftmark}
\fancyfoot[C]{\thepage}
\renewcommand{\headrulewidth}{0.4pt}

\setcounter{tocdepth}{2} % Keep ToC down to subsections

%everything before is preabmle, doc begins here

\begin{document}
\clearpage
\thispagestyle{empty}
\begin{center}
    \includegraphics[width=8cm,height=8cm,keepaspectratio]{UZH.png}
    \label{fig:uzh}
    
   	\vspace{2cm}
	\Huge
	\textbf{Efficient neural network wave functions for interacting fermions}
	
	\vspace{2cm}
	\huge
	Bachelor Thesis in Physics
	
	\vspace{2cm}	
	\Large
	\textbf{Illia Yakushevskyi}
	
	\vspace{3cm}
	\normalsize	
	Supervised by \\
	\Large
    Prof. Dr. Titus Neupert\\ 
	Dr. Attila Szabó\\
	Ph.D. candidate Simon Hille\\
	
	\vspace{2cm}
	\large
	\today
\end{center}
\vspace{\stretch{3}}
\clearpage
\pagenumbering{roman}
\thispagestyle{plain}
\vspace*{\stretch{1}}
\begin{abstract}

In the last few years, neural network wave functions have proven effective in describing many-body quantum systems. In particular, architectures such as FermiNet, PauliNet, and PsiFormer
have been successful in tackling systems of interacting electrons in quantum chemistry. However,
these all rely on representing antisymmetric fermion wave functions as determinants, whose evaluation may be prohibitively expensive for large systems.
A recently proposed alternative approach implements fermionic antisymmetry using a permutation symmetric neural network extended with a predefined antisymmetric function inspired by fractional quantum Hall states. In principle, such a neural network may be computationally more efficient than the standard approach.

\end{abstract}
\vspace{\stretch{3}}
\clearpage
\thispagestyle{plain}
\vspace*{\stretch{1}}
\begin{center}
   \textbf{Acknowledgements}
   \begin{quotation}
   XXX
   \end{quotation}
\end{center}
\vspace{\stretch{3}}
\clearpage
\newpage
\thispagestyle{plain}
\tableofcontents
\newpage
\listoffigures
\listoftables
\thispagestyle{plain}
\newpage
\stepcounter{page}
\pagenumbering{arabic}
\thispagestyle{plain}
\newpage
\section{Introduction}

In theoretical condensed matter in order to describe properties of material it is usually sufficient to describe the ground state of the many body system. However, it had always been a hard task to solve, which was mainly accomplished by various variational methods. It is becoming particularly hard to solve Quantum Many Body system, which include non-trivial strongly correlated interactions. 
Recently, relatively new approaches to model Wavefunction had emerged, such as Neural Quantum States (NQS)\cite{carleo_solving_2017}. Using it as the ansatz with the finite basis set in Variational Monte Carlo (VMC) techniques had improved the accuracy of the method significantly. 



While the computational cost of a bosonic ansatz can scale favorably, encoding the required spin statistics for fermions introduces a major structural challenge. For bosons, one can achieve permutation invariance through simple symmetric pooling, but fermions require a totally antisymmetric wavefunction to satisfy the Pauli Exclusion Principle. This necessitates the use of the Slater Determinant, using LU-decomposition it can be calculated in $O(n^3)$, which can be a bottleneck for training systems with large number of fermions. 

Different methods were proposed over the year to approach it, such calculating Slater Determinant, Pfafian ,  as backflow transformations initialy proposed by R.Feynman  \cite{feynman_energy_nodate} , ...  and later realised as neural backflow \cite{pfau_ab_2020}


.. Universal Approximation Theory tells us that with sufficiently large linear layer and non-linear activation function we can approximate function of arbitrary complexion, which is the premise of my thesis. 

% epxlain motivation for making separation of parts 

\newpage

% \setcounter{section}{1}



\thispagestyle{plain}


\section{Theoretical Background}

\subsection{Quantum Many-Body Problem in First Quantisation}



Fermions are the identical particles which.

Slater determinant 

\cite{orus_practical_2014} For a system of e.g. N spins 1/2, the dimension of the Hilbert space is 2N, which
is exponentially large in the number of particles. Therefore, representing a quantum state of the
many-body system just by giving the coefficients of the wave function in some local basis is an
inefficient representation. The Hilbert space of a quantum many-body system is a really big place
with an incredibly large number of quantum states. not all quantum states in the Hilbert space of a many-body system are equal: some are more relevant than others. 
In particular, one can prove that low-energy eigenstates of gapped Hamiltonians with local interactions obey the so-called area-law for the entanglement entropy [53], see Fig.(3). This means that the entanglement entropy of a region of space tends to scale, for large enough regions, as the size of the boundary of the region and not as the volume3. And this is a very remarkable property, because a quantum state picked at random from a many-body Hilbert space will most likely have a entanglement entropy between subregions that will scale like the volume, and not like the area. In other words, low-energy states of realistic Hamiltonians are not just “any” state in the Hilbert space: they are heavily constrained by locality so that they must obey the entanglement area-law.



\cite{lange_architectures_2024} 


+ The Fermionic Sign Problem

figure targetting part of hilbert space 

 formal mathematics of the permutation symmetric network extended with the predefined antisymmetric function. 

- Nqs works in contuniuous space 
Here some
\newpage

\subsection{Variational Monte Carlo}

 The base of this approach is Quantum Mechanical Variational Principle, one of the big advantages of which that the enery of the system is bounded from below by a ground state. By optimizing our variational state with respect to perameters $\theta$ in order to decrease energy we hope to capture the aproximation of the ground state. 


\subsubsection{Stochastic Reconfiguration Driver (SR)}


The big part of the VMC set up is a VMC driver. We use Netket's implementation of \texttt{VMC\_SR}  driver and this section will explain it. 

Generally the driver defines training loop which generates samples with Monte Carlo, evaluates the local energies $E_{loc}$ ,  calculates the gradients and updates parameters of $VMC_\theta$. 

The driver we use is called Stochastic Reconfiguration Driver with the SR preconditioner \cite{rende_simple_2024} \cite{becca_quantum_2017}. The role of SR is to  precondition the gradient used in stochastic optimization, improving the convergence rate to the ground state of a Hamiltonian by multiplying the gradients with the inverse of S matrix. 

S matrix is actually a Quantum Geometric Tensor (QGT). While standart ML optimizers like Adam work in euclidian parameter space and use Euclidian Norm, QGT is equiped with quantum mechanical distance function, Fubini-Study distance. 
To construct SR update we will first need gradient with respect to variational parameters $\theta$ 

We concerned with a problem to calculate optimal $VMC_\theta$. update and thus need to solve 

\begin{equation}
    \delta \theta =g S 
\end{equation}



Since S matrix is symmetric and positive semi-definite we use Cholesky decomposition for solving the inverse of the system. Here we need to introduce parameter we will use later which is diagonal shift. It halps with the ill-conditioned S matrixes. 


Explain it's parameters 


The bottleneck of fitting the model on a GPU is S matrix of Stochastic reconfiguration driver, which has dimmesions PxP, where P are the number of Parameters. 
Total number of parameters of the model can be calculated as: 
\begin{equation}
    P = 
\end{equation}

Show connection to Jacobian Matrix and how to calculate total weight


\newpage


\subsection{Neural Network Ansatz for Fermions}
\label{sec:nn_ansatz}

While for bosons -- for fermions Slater scales as ... 

for fermions efficiently representing antisymmetric states remains a major challenge. The standard aproach for fermions The standard approach is to expand $\Psi$ in terms of Slater determinants constructed from a single-particle basis set $ /phi (r), ..., φ_L(r).$
The problem with Slater determinant, is that it scales exponentially with number of N particles. From perspective of computational cost, it translates to the $O(n^3)$ complexity, when solving the Slater determinant with LU decomposition. The main goal of this thesis is to verify Ansatz alternative to the Slater determinant, which could bring down the computational cost to $O(n^2)$, thus beeing more efficient representation for fermionic wavefunctions. Such Ansatz was first proposed to use a symmetrical part and antisymmetric part as arguments to neural network. \cite{fu_minimal_2025}

\begin{equation}
\psi(r_1 ,..., r_n) = \Psi( r_1, ...,r_n,  \eta(r_1 ,..., r_n  ))
\end{equation}


Where $\Psi$ is implemented as Neural Network. While we know how to make symmetrized Ansatz for bosonic case \ref{tab:deepsets} the challenge is to construct so called \textit{signature encoder}:

\begin{equation}
    \eta = \prod_{1 \le j < i \le N} (x_i - x_j)
\end{equation}

Which can be extended to 2 dimmensions using the complex coordinates. 

\begin{equation}
    \eta = \eta_1 + i \eta_2 = \prod_{1 \le j < i \le N} (z_i - z_j)
\end{equation}


Later in this work we use regularised version of the $\eta$ for stabitily at the nodes when

$ |z_i - z_j| \rightarrow 0 $

\begin{equation}
\title{eq:eta}
\eta = \prod_{1 \le j < i \le N} \frac{z_i - z_j}{\sqrt{|z_i - z_j|^2 + a^2}}
\end{equation}

\title{tab:deepsets}
\subsubsection{Deep Sets}


... 

Insert the model graph 

\newpage



\section{Implementation and Experimental Setup}

System (Hamiltonian) $\rightarrow$ Sampler ( Metropolis Algo ) $\rightarrow$ Ansatz (NQS) $\rightarrow$ VMC (Eigen energies)

\subsection{System Architecture}

\subsection{The Sampler}

// explain what metrics we use for looking at the problem , Rhat, TauCorr etc 

\subsection{System}
\label{impl:sys}

// mb remove completely 

The object \textbf{System} includes description of continuous Hilbert Space we work with, it's dimension and Hamiltonian (namely the Potential part) we solve Schrödinger's equation for. Potential of our choice to begin with is Quantum Harmonic Oscillator. 
It's worth noting that we enforce antisymmetrization of the state into our Ansatz (Neural Quantum State) 

The configuration file of this experiment \texttt{Experiment}\footnote{\url{https://github.com/IlyaYakushevskiy/FermiNQS/blob/main/configs/experiment/qho_fermisets_2d_5N_post.yaml}}

\texttt{Ansatz}\footnote{\url{https://github.com/IlyaYakushevskiy/FermiNQS/blob/main/src/ansatz.py}} .



Simulations for this work were performed using 
NetKet \cite{carleo_netket_2019, vicentini_netket_2022}. This software is built on top of JAX~\cite{bradbury_jax_2018} 
and Flax~\cite{heek_flax_2024}.

\subsection{Neural Network Ansatz}

As was introduced in section \ref{sec:nn_ansatz} ...

The size of our input layer is defined by several parameters. 

- to avoid correlation between the Markov Chains we increase \texttt{n\_chains}, which determines amount of chains we spawn 

- \texttt{n\_samples} defines amount of samples generated by each chain

- \texttt{N} is amount of particles in the symtem

- \texttt{dim} is respectively a dimension of continuous Hilbert space 

- additional computational parameter is \texttt{chunk\_size} which allows us to fit the model on GPU. 

$ input_dimmension = n\_chains *  $

another variables in the ansatz are : 
- \texttt{a} is a regularisation parameter for antisymmetric function. 

Another hurdle of computational implementation is that we always have to treat output of our Ansatz as logarithm.  


 


\newpage

Some examples
\begin{equation}
    M_Z=(57.9\pm7.03)M_{\oplus} (\frac{M}{M_J})^{(0.61\pm 0.08)}
\end{equation}
\begin{table}[H]
    \centering
    \begin{tabular}{|c||c|c|c|c|c|}
        \hline
         n&0&1&1.5&2&2.5\\
         \hline
         $\tau_{ff}$&$\propto M^0$&$\propto M^{1/4}$&$\propto M^{1/2}$&$\propto M$&$\propto M^{5/2}$\\
          $\tau_{expl}$&$\propto M^0$&$\propto M^{1/4}$&$\propto M^{1/2}$&$\propto M$&$\propto M^{5/2}$\\
          $\tau_{Alfven}$&$\propto M^{1/3}$&$\propto M^{1/4}$&$\propto M^{1/6}$&$\propto M^0$&$\propto M^{-1/2}$\\
         $\tau_{shear}$&$\propto M^{-1/3}$&$\propto M^{0}$&$\propto M^{1/3}$&$\propto M$&$\propto M^{3}$\\
         \hline
    \end{tabular}
    \caption[Short Caption Table 1]{\textit{Relation of the timescales with planetary mass for different density power-law indexes.}}
    \label{tab:n}
\end{table}
\begin{figure}[H]
     \centering
     \begin{subfigure}[b]{0.495\linewidth}
         \centering
         \includegraphics[width=\linewidth]{cormassa.png}
         \caption{\textbf{free-fall timescale}}
         \label{fig:ratio3}
     \end{subfigure}
     \begin{subfigure}[b]{0.495\linewidth}
         \centering
         \includegraphics[width=\linewidth]{cormassb.png}
         \caption{\textbf{expansion timescale}}
         \label{fig:ratio4}
     \end{subfigure}
     \begin{subfigure}[b]{0.495\linewidth}
         \centering
         \includegraphics[width=\linewidth]{cormassc.png}
         \caption{\textbf{Alfvén timescale}}
         \label{fig:ratio5}
     \end{subfigure}
     \begin{subfigure}[b]{0.495\linewidth}
         \centering
         \includegraphics[width=\linewidth]{cormassd.png}
         \caption{\textbf{shear timescale}}
         \label{fig:ratio6}
     \end{subfigure}
     \caption[Short Caption Figure 1]{\textit{Relation of the 13 MHD clumps bound mass with (a) the free-fall timescale, (b) the expansion timescale, (c) the Alfvén timescale, and (d) the shear timescale. The black dashed and dotted lines correspond to the predicted relation for power-law index n=1.5 and n=2.}}
\end{figure}
\begin{table}[H]
    \centering
    \begin{tabular}{|l||r|r|r|r|}
    \hline
     Mass & $0.01M_{Jupiter}$ & $0.1M_{Jupiter}$& $1M_{Jupiter}$& $10M_{Jupiter}$\\
     \hline
    \hline
    \multicolumn{5}{|c|}{ \textbf{Vaporization timescale} [$10^3$yrs]}\\
    \hline
    \hline
   Z=0.007 &500'000&5'000&50&0.5\\
   Z=0.013 &900'000&9'000&90&0.9\\
   Z=0.02 & 1'300'000&13'000&130&1.3\\
   Z=0.1 & 4'300'000&43'000&430&4.3\\
   Z=0.2 & 4'800'000&48'000&480&4.8\\
   Z=0.5 & 1'100'000&11'000&110&1.1\\
   Z=0.7 & 130'000&1'300&13&0.1\\
   Z=0.9 & 660&6.6&0.07& 0.001\\
     \hline
\end{tabular}
 \caption[Short Caption Table 2]{\textit{Vaporization timescales of a micron-sized silicate grain in scenario 1. } }
 \label{tab:vap}
\end{table}
Figure \ref{fig:planets} shows...


\newpage

\section{Results $\&$ Discussion}

For the system discussed in \ref{impl:sys} we going to do an ablation studies according to \cite{godbole_deep_2023}. 


We going to start with a simplest system, 4 spinless fermions in Quantum Harmonic oscilator potential with the following Hamiltonian to test our ansatz : 

\begin{equation}
\hat{H}_0
=
\sum_{i=1}^{4}
\left[
    -\frac{\hbar^{2}}{2m}\nabla_i^{2}
    + \frac{1}{2} m \omega^{2} \mathbf r_i^{2}
\right]
\end{equation}

Let us look first at 2D case, we define $ m  \omega^2 = 1  $ , so we can find analytic ground state as $ \hbar \omega ( n + 1 )  $. The problem with training the network for 4 fermion arises from the three-degenerate ground state. 


From analysis of possible $n = n_x + n_y$ single particle wavefunction we can see closest states to groundstate will be : 




(n = 0) : (0,0)

(n = 1) : (0,1), (1,0)

(n = 2): (0,2), (2,0), (1,1) 

(n = 3) : (1,2), (2,1), (3,0), (0,3)

 
So the true Ground state would be of the form: 

\begin{equation}
\Psi_{\mathrm{GS}}
=
c_1\,\Psi_{2,0}
+
c_2\,\Psi_{1,1}
+
c_3\,\Psi_{0,2}
\end{equation}

and the Energy of this system would be : 
\begin{equation}
    E = \hbar \omega ( n_x + n_y + \frac{1}{2} + \frac{1}{2}) =  \hbar \omega  (( 0 + 0 + 1 ) + ( 0 + 1 + 1  )  +  ( 1 + 0 + 1  ) + ( 0 + 2 + 1  )) =  \hbar \omega 8
\end{equation}


From our energy convergece graph \ref{fig:n4energy} we see that all configurations are heavily biased to the state with E = 10 , which, as testing with different N confirms is always the exact Laughlin state. 

\subsection{Test of Optimizers}

Since we had main troubles with instabilites within SR driver, we also tested the very standard general Machine Learning aproach to use Adam Optimizer without stochastic reconfiguration. Since we don't have to compute S matrix that allowed us to fit much bigger model on the same GPU. 
Instead of having [16,32,64] hidden units we used  [16,32,1024] , and also additional configuration we added 1 extra layer with weight size of dimmension $[hidden_size x hidden_size]$

We report that adding extra layer (even decreasing ) increased training time massively (30 times more per iteration), which due to the way code is paralelized in our set up. 

Set up with Adam and no SR did not have numerical instability troubles and found Laughlin state with low errors around 0.001 percent, but couldn't converge to the actual ground state  

\paragraph{Solver for SR Optimizer}


However, the inverse is often ill defined.  One reason is that Monte Carlo estimates of matrix elements are noisy. Noise accumulates to render the matrix singular by making a small eigenvalues vanish. Therefore, quickly and efficiently obtaining many uncorrelated samples from is crucial. \cite{medvidovic_variational_2023}

\subsection{Hyperparameter Optimization}
\paragraph{Effect of hidden layer size}

\begin{figure}
    \centering
    \includegraphics[width=15cm,height=10cm,keepaspectratio]{plots/Mean_energy_convergence_for_N=4_and_varied_hidden_size.png}
     \caption[Short Caption Figure 2]{\textit{My caption}}
     \label{fig:n4energy}
\end{figure}


best \texttt{diag\_shift} is 0.001 with soft SVD optimiser. 
\paragraph{Regularisation of antisymmetric part}

In our formualation of antisymmetric function \ref{eq:eta}
it's regularised by parameter a, we a made a quick sweep to assses it's effect: 


\subsection{Performance of the best model}

best result is run 2026-05-14/18-38-51

graph for variance -- refer to zero varience principle
Overlap over time - oscinlating behaviour ? 
TODO plot agains laughlin state

main problem is characterised by smooth descent -> jump -> smooth descent

an idea for it :  
"Continuous functions cannot just teleport their zero-crossings. To fix the geometry, the network has to literally tear and reconnect these multidimensional nodal surfaces.

During the steps where the nodes are actively shifting and crossing through regions of high particle density, the kinetic energy violently spikes. Once the nodes snap into the correct new topology, the energy collapses back down and resumes a smooth descent toward a lower plateau."
-> read on mitigation strategies 
    - energy clipping 
    - metropolis exchange 

print MSCATATE PARAMETERS 

LIST THE PLOTS 

EXPORT OUTPUTS / folder ON HDD 

INCLUDE 1D EXPL / COMPTA TO 2D



\subsection{Laughlin State (MAIN MEAT3)}

Classically the Lowest Landau level of Laughlin state is  $\Psi_{LL} = Vandermonde \cdot Boson \cdot Gaussuian  $
and it's a state constructed to represent electron gas placed in uniform background magnetic field in the presence of a uniform jellium background.  

Empirically it also appears to be "lazy" state which our ansatz \ref{} conveges to. //go back to L Fu and right down the ansatz 

Proof: 

Theory: It gets stuck in Laughlin State () , compute overlap with laughlin state, compute overlap between (state, state+ 50 iters) 

\subsection{Supervised pretraining}

As we seen in previous section the shape of our Ansatz leads VMC to get stuck in Laughlin eigenstate trap. The natural question was if expressivity of the Ansatz was a problem. We thus made supervised pretraing for our 2D QHO setup and in the post training the eigen-energy converged to the energy of real ground state. This rules out the idea that Ansatz lacks expressivity and can't replicate the structure of Slater determinant.
<paste the plot>

\subsection{Results with radial symmetry assumption}

Where finally achieved first success was in introducing the radial symmetry for 2D QHO, this is a special case which works when the ground state is  radially symmetrical, this means we must have spherical potential as well as GS must be non-degenerate. 

We thus can work with systems of 3,6,10,... particles. Later we will test it on quantum dot system, where we encode the Coulomb interaction \ref{sec:q_dot} 

For such system we know for fact that $ [\hat{H}, \hat{L_z}] = 0$, meaning that Eigen energies of the Hamiltonian are also eigenenergies of $\hat{L_z}$ operator. 

We thus need a projection of our neural network onto $\hat{L_z}$. The way we achieve it is to apply rotation matrix k times with angles $ \theta = \frac{k}{2 \Pi}$  to the coordinates input into symmetrical part (Nerual Network) and then average it out. We must note that this increases the computational cost by factor of k. Hence the cost of ansatz is $O(k \cdot n^2 )$ compared to $O( n^3)$ in case of the Slater Determinant ansatz. The scaling of error with k and N is yet to be studied. 




\newpage

\section{Conclusions $\&$ Outlook}
\subsection*{Conclusions}
\subsection*{Outlook}
\newpage
\part{Appendix}
\bibliographystyle{ieeetr}
\bibliography{zotero_references}
%\printbibliography
\end{document}
